\subsection{Revue des objectifs}

Ce laboratoire a permis d'étudier expérimentalement les ondes à travers deux manipulations, l'étude des propriétés des micro-ondes à l'aide d'un émetteur à cornet, puis la mesure de la vitesse du son dans les gaz à l'aide du trombone de Koenig, afin de confronter les mesures aux modèles théoriques vus au cours.

\textbf{\textit{Manipulation 1 : Ondes centimétriques -- Micro-ondes}}

\textbf{Objectif 1 : Propagation rectiligne}

Le signal a été mesuré en déplaçant le récepteur perpendiculairement à l'axe du faisceau et il est maximal lorsque les cornets sont alignés. L'objectif est atteint.

\textbf{Objectif 2 : Capacité de pénétration}

La tension a été mesurée derrière une plaque de bois à plusieurs distances et le bois laisse passer une partie du signal. L'objectif est atteint qualitativement.

\textbf{Objectif 3 : Blindage et absorption}

La tension a été mesurée sans plaque et derrière une plaque d'aluminium, et aucun signal n'est détecté derrière l'aluminium. L'objectif est atteint.

\textbf{Objectif 4 : Réflexion}

L'angle de réflexion a été mesuré pour plusieurs angles de plaque et les écarts avec la loi de la réflexion sont inférieurs aux incertitudes. L'objectif est atteint quantitativement.

\textbf{Objectif 5 : Ondes stationnaires}

La longueur d'onde a été déterminée à partir d'une seule paire minimum-maximum, ce qui rend la mesure peu précise. L'objectif n'est atteint que partiellement.

\textbf{Objectif 6 : Réfraction}

La déviation d'un prisme de paraffine a été mesurée une seule fois et l'indice obtenu est nettement inférieur à la valeur attendue. L'objectif n'est atteint que partiellement, la méthode de la déviation minimale s'appliquant mal à ce montage.

\textbf{Objectif 7 : Principe de Huygens}

La tension dans la zone d'ombre a été mesurée avec et sans plaque à fente, et elle augmente nettement avec la fente. L'objectif est atteint qualitativement.

\textbf{Objectif 8 : Diffraction}

Le signal a été mesuré derrière une plaque couvrante et il subsiste dans la zone d'ombre en augmentant vers les bords. L'objectif est atteint qualitativement.

\textbf{Objectif 9 : Interférence}

Le signal a été mesuré derrière une double fente et une alternance de maxima et de minima est observée, bien que le pas de mesure soit trop grand pour bien les résoudre. L'objectif est atteint qualitativement.

\textbf{Objectif 10 : Polarisation}

Le signal a été mesuré pour différentes orientations de la grille de polarisation et il passe d'un maximum à une valeur quasi nulle selon l'orientation des fentes. L'objectif est atteint.

\medskip

\textbf{\textit{Manipulation 2 : Trombone de Koenig -- Vitesse du son dans les gaz}}

\textbf{Objectif 1 : Vitesse théorique}

La vitesse théorique du son a été calculée avec son incertitude pour l'air et pour le CO$_2$ à partir des températures mesurées. L'objectif est atteint.

\textbf{Objectif 2 : Mesure dans l'air}

La vitesse du son dans l'air a été mesurée à plusieurs fréquences et déterminée par une régression linéaire passant par l'origine. L'objectif est atteint.

\textbf{Objectif 3 : Mesure dans le CO$_2$}

La vitesse du son dans le CO$_2$ a été mesurée à plusieurs fréquences et déterminée par une régression linéaire passant par l'origine. L'objectif est atteint.

\textbf{Objectif 4 : Comparaison}

Pour les deux gaz, la vitesse est indépendante de la fréquence et les écarts avec les valeurs théoriques restent faibles, bien que supérieurs aux incertitudes. L'objectif est atteint.

Les mesures dans l'hélium, proposées par le descriptif, n'ont pas été réalisées car le gaz n'était pas disponible.
